\SourceLeaf{44}
\fancyhead[L]{\small DETERMINING THE ELEVATION OF THE POLE}
\renewcommand{\Article}[1]{\par\medskip\Needspace{5\baselineskip}\pdfbookmark[1]{Art. #1}{poli1808-art-#1}\begin{center}\textbf{#1.}\end{center}\nobreak}
\pdfbookmark[0]{A Particular Method for Determining the Elevation of the Pole}{poli1808-body}
\begin{center}\large A PARTICULAR METHOD\par\medskip
FOR DETERMINING THE ELEVATION OF THE POLE.\par\medskip\rule{.09\linewidth}{.3pt}\end{center}
\Article{1}
The methods most commonly used by traveling observers to determine geographical position are known to rely on observations of the Sun. Corresponding altitudes serve to determine the time, and altitudes at or near the meridian serve to determine the elevation of the pole. This kind of observation has several advantages, especially for those who measure altitudes with Hadley’s sextant. With this instrument it is much easier to observe the Sun’s altitude accurately than the altitudes of the stars, and the divisions on the graduated limb are harder to read by night than by day. Moreover, the calculation from corresponding altitudes and those near the meridian is so simple that anyone can readily learn it, even with little mathematical experience. --- If clouds or other obstacles make it impossible to obtain corresponding altitudes, observations close to culmination are usually combined with others farther from it. The difference between the clock and the true time, as well as the altitude of the pole, is generally derived from them by indirect methods, such as the one named after \textsc{Cornelis Douwes}.

Although there is no doubt that observations of the Sun are preferable to nighttime observations whenever they can be made, we think the latter deserve attention too. Often, after a cloudy day, the sky clears as night falls or at least allows occasional glimpses of
\SourceLeaf{45}the stars. Often, too, a traveler who had no time for observations during the day because he was intent on continuing his journey could render a service by determining the geographical position of the place where he spends the night, provided he has acquired enough skill in observing the stars. Add to this the uncertainty of observations of the Sun in the winter months, when in our regions the Sun rises scarcely above the haze on the horizon, and the more frequent interference from intrusive onlookers and the jolts that passing vehicles give an artificial horizon, especially a mercury horizon. In the quiet of the night you are better protected from such disturbances. We must admit that observations of the stars are well worth cultivating more carefully and putting to use.

The benefit this kind of observation can offer \textit{on voyages at sea} is evident when one considers how important it is to ascertain a ship’s position, daily if possible. People’s safety and lives often depend on knowing it. Nighttime observations at sea admittedly have difficulties of their own, since the sea horizon, from which the altitudes must be measured, is harder to make out at night. This difficulty can be remedied, however, by fitting the sextant with a telescope of wider aperture.
\Article{2}
The arguments in favor of stellar observations given in the preceding article carry even more weight when we consider that this method can determine both the altitude of the pole and the time with remarkable precision in a very short interval, for example from observations made within the first few minutes. For this purpose one need only observe the altitudes of any two different stars whose positions are assumed known and note the corresponding readings of the clock. The daily rate of the clock need be known only accurately enough to convert the time interval into sidereal time with sufficient precision. Unless we are mistaken, this problem surely ranks among the most useful in nautical astronomy, and it is rather surprising that those who have written on the subject seem to have neglected it entirely. The distinguished \textsc{Kraft} showed in \textit{Actis novis Acad. Petropol.}, vol.~XIII, how to determine the elevation of the pole from the altitudes of two stars. For a reason I do not know, he restricted the general problem by imposing a special condition, namely that the two observations were
\SourceLeaf{46}made at exactly the same instant. This calls for two observers, two instruments, and exact simultaneity of the observations, which is certainly rather difficult to achieve. Yet this condition is quite useless and superfluous. A well-equipped ship surely has a clock that can measure an interval of a few minutes accurately, and the ship’s motion over so short a time cannot produce an appreciable change in its geographical position. Finally, this condition makes the calculation for solving the problem no simpler, as we shall show at once.
\Article{3}
It is easy to see that the problem, in the form treated by the distinguished \textsc{Kraft}, can also be stated as follows. ``Given the positions of two points on a spherical surface relative to a great circle (here the equator), find the position of a third point whose distances from those points are given.'' The first two points are given by the positions of the two observed stars on the celestial sphere, for instance by their right ascensions and declinations. The third point corresponds to the zenith of the place of observation. Its declination determines the elevation of the pole, and its right ascension determines the culminating point of the equator and hence the sidereal time. The solution that suggests itself at once uses three spherical triangles. In the first, formed by the two given points and the pole of the equator, two sides and the included angle are given. From them we obtain the third side and either of the two remaining angles, for example the one at the first point. In the second triangle, formed by the two given points and the third point sought, all the sides are already given. From them we obtain one of the angles, namely the one again at the first point if we chose the angle at that point in the first triangle. The sum or difference of these two angles is the known angle in the third triangle, formed by the first and third points and the pole of the equator. The two sides adjacent to this angle are also known. Finally, from these we obtain the angle at the pole and the distance of the third point from the pole, which is the complement of the elevation of the pole. Astronomers were already using this method in the sixteenth century. See, for example, \textit{Tychonis Astronomiae instauratae progymnasmata}, p.~211 and following, where the positions of a great many stars, including the new star in Cassiopeia, are determined from their distances from pairs of stars of known position.

\SourceLeaf{47}\SourceGathering{VI.~6}It is also clear from this that, strictly speaking, the problem always admits two solutions. For if two circles are drawn on the spherical surface about the two given points as centers, with radii equal to the respective given distances, the circles will generally intersect at two points, and either intersection may be taken as the third point sought. In practice, however, this cannot give rise to ambiguity. Imagine the great circle drawn through the first and second points, dividing the spherical surface into two hemispheres. There can be no doubt whether the third point, the zenith itself, and the north pole must lie in the same hemisphere or in opposite ones. The first case must apply if the excess of the right ascension of the star that was on the left of the other during the observation over the right ascension of the other star, which was to its right, lies, with $360$ degrees added if necessary, between $0$ and $180^0$. The second case applies if this excess lies between $180^0$ and $360^0$.

The lengthy calculation that burdens the direct solution is due to this very ambiguity, which is inherent in the problem itself. In what follows, however, we shall show how to shorten it considerably. If we are willing to use an indirect solution, a fairly quick one can be devised. We reserve its explanation for another occasion.

So far we have discussed the special case in which the two observations are simultaneous. The general problem reduces to this case in the following way. In place of the star observed second, imagine a fictitious star with the same declination and a right ascension reduced by the sidereal time elapsed between the first and second observations. This fictitious star would clearly have reached at the time of the first observation the same altitude that the star actually observed had at the time of the second observation. Once the fictitious star has been substituted for the actual one, the calculation can be based on simultaneous observations.

All this was drawn from readily apparent but geometric considerations. We therefore hope that a direct solution based solely on analytic operations will be welcome to many readers. It will surely offer further confirmation that hardly anything can be obtained from geometric considerations that cannot be obtained just as neatly by analysis, if handled skillfully.
\SourceLeaf{48}
\begingroup
\newcommand{\GaussVIPoliLabel}[3]{%
  \makebox[0pt][l]{[#2]}%
  \makebox[#1\linewidth][r]{$\displaystyle #3$}%
}
\newcommand{\GaussVIPoliAligned}[1]{%
  \[\makebox[\linewidth][l]{$\displaystyle\begin{aligned}#1\end{aligned}$}\]%
}
\newcommand{\GaussVIPoliSingle}[2]{%
  \[\makebox[\linewidth][l]{%
    \makebox[0pt][l]{[#1]}%
    \makebox[\linewidth][c]{$\displaystyle #2$}}\]%
}
\Article{4}
Let $\varphi$ be the elevation of the pole, $\alpha$ and $\alpha'$ the right ascensions, $\delta$ and $\delta'$ the declinations of the two stars, and $\gamma$ and $\gamma'$ the right ascensions of the culminating points of the equator at the times of the first and second observations, respectively, or equivalently the sidereal times themselves converted to degrees. Finally, let $h$ be the observed altitude of the first star and $h'$ the altitude of the second. Evidently $\gamma-\alpha$ and $\gamma'-\alpha'$ are the hour angles corresponding to the two observations. We shall set $\gamma-\alpha=\lambda$ and $\gamma'-\alpha'=\lambda-\theta$. Then $\theta=\alpha'-\alpha-(\gamma'-\gamma)$ is known, since $\gamma'-\gamma$ equals the interval between the two observations converted to sidereal time and then to degrees. The solution of the problem is therefore to be based on these two equations:
\GaussVIPoliAligned{
\GaussVIPoliLabel{.30}{1}{\sin h}&=\sin\delta\sin\varphi+\cos\delta\cos\varphi\cos\lambda\\
\GaussVIPoliLabel{.30}{2}{\sin h'}&=\sin\delta'\sin\varphi+\cos\delta'\cos\varphi\cos(\lambda-\theta)
}
If we tried to eliminate either the unknown $\lambda$ or the unknown $\varphi$ from these equations, we would arrive at an overly complicated equation. The distinguished \textsc{Kraft} took the former course in the paper cited above. In our judgment, however, the solution he obtained this way is even longer and more laborious than the one derived directly from the three triangles, even though his solution determines only $\varphi$ and leaves the time undetermined. It is therefore better to introduce a new unknown and use it to eliminate both $\varphi,\lambda$. The following observation will guide us in choosing it. We have
\[
\begin{aligned}
&(\sin\delta\sin\varphi+\cos\delta\cos\varphi\cos\lambda)^2+(\cos\delta\sin\varphi-\sin\delta\cos\varphi\cos\lambda)^2\\
&\qquad=\sin\varphi^2+\cos\varphi^2\cos\lambda^2=1-\cos\varphi^2\sin\lambda^2
\end{aligned}
\]
and hence by equation $1$
\[
(\cos\delta\sin\varphi-\sin\delta\cos\varphi\cos\lambda)^2+\cos\varphi^2\sin\lambda^2=\cos h^2,
\]
or
\[
\left(\frac{\cos\delta\sin\varphi-\sin\delta\cos\varphi\cos\lambda}{\cos h}\right)^2+\left(\frac{\cos\varphi\sin\lambda}{\cos h}\right)^2=1
\]

We may therefore set
\GaussVIPoliAligned{
\GaussVIPoliLabel{.56}{3}{\dfrac{\cos\delta\sin\varphi-\sin\delta\cos\varphi\cos\lambda}{\cos h}}&=\cos u\\[4pt]
\GaussVIPoliLabel{.56}{4}{\dfrac{\cos\varphi\sin\lambda}{\cos h}}&=\sin u
}

\SourceLeaf{49}\SourceGathering{6\textsuperscript{*}}Now, combining equation $3$ with $1$ gives
\GaussVIPoliAligned{
\GaussVIPoliLabel{.34}{5}{\sin\varphi}&=\sin\delta\sin h+\cos\delta\cos h\cos u\\
\GaussVIPoliLabel{.34}{6}{\cos\varphi\cos\lambda}&=\cos\delta\sin h-\sin\delta\cos h\cos u
}

Let us write equation $2$ in the form
\[
\sin h'=\sin\delta'\sin\varphi+\cos\theta\cos\delta'\cos\varphi\cos\lambda+\sin\theta\cos\delta'\cos\varphi\sin\lambda
\]
and substitute for $\sin\varphi$, $\cos\varphi\cos\lambda$, $\cos\varphi\sin\lambda$ their values from equations $5,\ 6,\ 4$, which gives
\[
\begin{aligned}
&\sin h'-\sin h\sin\delta\sin\delta'-\sin h\cos\theta\cos\delta\cos\delta'-\cos u\cos h\cos\delta\sin\delta'\\
&\quad+\cos u\cos h\cos\theta\sin\delta\cos\delta'-\sin u\cos h\sin\theta\cos\delta'=0
\end{aligned}
\]

If we set
\GaussVIPoliSingle{7}{\frac{\cos\delta\sin\delta'-\cos\theta\sin\delta\cos\delta'}{\sin\theta\cos\delta'}=\operatorname{cotang}v}
this equation becomes
\[
\begin{aligned}
&\sin h'-\sin h\sin\delta\sin\delta'-\sin h\cos\theta\cos\delta\cos\delta'\\
&\qquad-\cos h\sin\theta\cos\delta'(\cos u\operatorname{cotang}v+\sin u)=0
\end{aligned}
\]

So if we set $v-u=w$, we have
\GaussVIPoliSingle{8}{\cos w=\frac{\sin v(\sin h'-\sin h\sin\delta\sin\delta'-\sin h\cos\theta\cos\delta\cos\delta')}{\cos h\sin\theta\cos\delta'}}
which determines $w$ and hence $u=v-w$.

Once the angle $u$ has been found, $\lambda$ is obtained by combining equations $4$ and $6$, which gives
\GaussVIPoliSingle{9}{\operatorname{tang}\lambda=\frac{\cos h\sin u}{\cos\delta\sin h-\sin\delta\cos h\cos u}}

From $\lambda$ we obtain $\gamma=\alpha+\lambda$. Converting this angle to time gives the clock error. Finally, combining equations $4$ and $5$ gives
\GaussVIPoliSingle{10}{\operatorname{tang}\varphi=\frac{\sin\lambda\,(\sin\delta\sin h+\cos\delta\cos h\cos u)}{\cos h\sin u}}
If desired, equation $4$ itself can be used alongside this formula to check the calculation.
\SourceLeaf{50}\Article{5}
We add a few observations about the solution given in the preceding article.

Formula $8$ gives two values for the angle $w$, as determined by its cosine, one positive and the other negative. Clearly this gives two different solutions, as we already noted above. It is easy, however, to tell which one is the true solution in the following way. It can be shown that the sine of the difference between the azimuth of the vertical circle in which the first star was observed and that of the vertical circle in which the second star was observed, or more precisely the sine of the excess of the former azimuth over the latter, measured from left to right, is
\[
=\frac{\sin\theta\cos\delta'\sin w}{\cos h'\sin v}
\]
Now since $\cos\delta',\ \cos h'$ are by their nature positive quantities, clearly $\sin w$ must have either the same sign as $\frac{\sin\theta}{\sin v}$ or the opposite sign, according as the first vertical circle lay to the right or to the left of the second. There can be no uncertainty about this, since observations made in vertical circles that nearly coincide or are nearly opposite are quite unsuitable in practice and must therefore be carefully avoided, as we shall explain more fully below.

The solution can be shortened a little further and made more suitable for numerical calculation by introducing two auxiliary angles. In fact, if the angle $F$ is determined by the equation
\GaussVIPoliSingle{11}{\operatorname{tang}F=\frac{\operatorname{tang}\delta'}{\cos\theta}}
equations $7$ and $8$ become
\GaussVIPoliAligned{
\GaussVIPoliLabel{.34}{12}{\operatorname{tang}v}&=\dfrac{\cos F\operatorname{tang}\theta}{\sin(F-\delta)}\\[5pt]
\GaussVIPoliLabel{.34}{13}{\cos w}&=\dfrac{\cos v\operatorname{tang}h}{\operatorname{tang}(F-\delta)}\left(\dfrac{\sin h'\sin F}{\sin h\sin\delta'\cos(F-\delta)}-1\right)
}

Likewise, if we introduce an auxiliary angle $G$ such that
\GaussVIPoliSingle{14}{\operatorname{tang}G=\frac{\operatorname{tang}h}{\cos u}}
equations $9$ and $10$ will become
\GaussVIPoliAligned{
\GaussVIPoliLabel{.34}{15}{\operatorname{tang}\lambda}&=\dfrac{\cos G\operatorname{tang}u}{\sin(G-\delta)}\\[5pt]
\GaussVIPoliLabel{.34}{16}{\operatorname{tang}\varphi}&=\cos\lambda\operatorname{cotang}(G-\delta)
}
\SourceLeaf{51}\Article{6}
Practical astronomers often complain that methods proposed by theorists are sometimes unsuitable for practical use, or at least fall short of what they expected from them. To dispel any such suspicion about our method and at the same time illustrate how to use the preceding instructions, it will be useful to work through an example from actual observations in some detail. On August $21$ of this year, at our observatory, I observed the uncorrected double altitude of the star $\alpha$ Aquilae with a Troughton sextant with a radius of ten inches. It was $91^0\ 34'\ 10''$ when the Shelton clock showed $20^{\mathrm h}\ 40^{\mathrm m}\ 8^{\mathrm s}$. The star $\alpha$ Andromedae reached the same altitude at the clock reading $20^{\mathrm h}\ 46^{\mathrm m}\ 59^{\mathrm s}$. The first star was observed in the vertical circle to the right, the second in the one to the left. The daily rate of the clock agreed exactly with true sidereal time. Removing the index error from the observed angle gives the true angle $91^0\ 31'\ 36''$ and thus the apparent altitude $45^0\ 45'\ 48''$, from which we obtain the true altitude $45^0\ 44'\ 52''6=h=h'$. With aberration and nutation duly applied, the positions of the stars were found to be as follows:
\[
\begin{array}{rcl@{\quad}rcl}
\alpha&=&295^0\ 22'\ 6''6&\delta&=&+8^0\ 22'\ 43''1\\
\alpha'&=&359\ 38\ 18.5&\delta'&=&+28\ 2\ 13.4
\end{array}
\]
The time difference $6'\ 51''$ converted to an arc gives $\gamma'-\gamma=1^0\ 42'\ 45''$, whence $\theta=62^0\ 33'\ 26''9$. The calculation itself then runs as follows (the factor $\frac{\sin h'}{\sin h}$ in formula $13$ clearly drops out here because the altitudes are equal).
\[
\renewcommand{\arraystretch}{.87}
\begin{array}{@{}l@{\quad}r@{}l@{\hspace{.65em}}l@{}}
\log\operatorname{tang}\delta'\quad .\ .\ .\ .&9.7263516&&\\
\log\cos\theta\quad .\ .\ .\ .\ .&9.6635677&&\\\cline{1-3}
\log\operatorname{tang}F\quad .\ .\ .\ .&0.0627839&,&\text{whence }F=49^0\ 7'\ 37''70\\
&&&\text{and }F-\delta=40^0\ 44'\ 54''60\\[6pt]
\log\operatorname{tang}\theta.\quad .\ .\ .\ .&0.2845877&&\\
\log\cos F\quad .\ .\ .\ .\ .&9.8158318&&\\
\text{Comp. }\log\sin(F-\delta)\quad .&0.1802598&&\\\cline{1-3}
\log\operatorname{tang}v\quad .\ .\ .\ .\ .&0.2856793&&\text{whence }v=62^0\ 36'\ 58''79\\[6pt]
\log\sin F\quad .\ .\ .\ .\ .&9.8786157&&\\
\text{Comp. }\log\sin\delta'\quad .\ .\ .&0.3278629&&\\
\text{Comp. }\log\cos(F-\delta)\quad .&0.1205703&&\\\cline{1-3}
&0.3270489&&=\log 2.123483=\log n
\end{array}
\]
\endgroup
